% Geometric definition of the Euler angles (passive z-x-z convention, as in the text).
% All body axes are computed from (phi, theta, psi) with R = R_z(psi) R_x(theta) R_z(phi):
% the rows of R are the body axes x1, x2, x3 expressed in the inertial frame (X, Y, Z).
% Compile: xelatex euler_angles_tikz.tex ; pdftoppm -png -r 300 -singlefile euler_angles_tikz.pdf euler_angles
\documentclass[tikz,border=4pt]{standalone}
\usepackage{amsmath}
\usepackage[fontset=fandol]{ctex}
\usepackage{tikz-3dplot}
\usetikzlibrary{arrows.meta}

\begin{document}
\tdplotsetmaincoords{66}{125}
\begin{tikzpicture}[tdplot_main_coords,>=Stealth,line cap=round,
                    every node/.style={font=\normalsize}]
  % ---- Euler angles (degrees) used for the drawing
  \pgfmathsetmacro{\ph}{50}
  \pgfmathsetmacro{\th}{38}
  \pgfmathsetmacro{\ps}{40}
  % ---- line of nodes ON and the in-plane direction e' (body x1x2-plane, perpendicular to ON)
  \pgfmathsetmacro{\nx}{cos(\ph)}            \pgfmathsetmacro{\ny}{sin(\ph)}
  \pgfmathsetmacro{\ex}{-cos(\th)*sin(\ph)}  \pgfmathsetmacro{\ey}{cos(\th)*cos(\ph)}  \pgfmathsetmacro{\ez}{sin(\th)}
  % ---- body axes = rows of R
  \pgfmathsetmacro{\ax}{cos(\ps)*\nx+sin(\ps)*\ex}  \pgfmathsetmacro{\ay}{cos(\ps)*\ny+sin(\ps)*\ey}  \pgfmathsetmacro{\az}{sin(\ps)*\ez}
  \pgfmathsetmacro{\bx}{-sin(\ps)*\nx+cos(\ps)*\ex} \pgfmathsetmacro{\by}{-sin(\ps)*\ny+cos(\ps)*\ey} \pgfmathsetmacro{\bz}{cos(\ps)*\ez}
  \pgfmathsetmacro{\cx}{sin(\th)*sin(\ph)}  \pgfmathsetmacro{\cy}{-sin(\th)*cos(\ph)}  \pgfmathsetmacro{\cz}{cos(\th)}
  \def\L{3.0}

  % ---- inertial XY plane (disc) and body x1x2 plane (disc)
  \fill[gray!12] plot[domain=0:360,samples=90,variable=\t] ({2.7*cos(\t)},{2.7*sin(\t)},0) -- cycle;
  \draw[gray!50] plot[domain=0:360,samples=90,variable=\t] ({2.7*cos(\t)},{2.7*sin(\t)},0) -- cycle;
  \fill[blue!10,opacity=.6] plot[domain=0:360,samples=90,variable=\t]
      ({2.3*(cos(\t)*\nx+sin(\t)*\ex)},{2.3*(cos(\t)*\ny+sin(\t)*\ey)},{2.3*sin(\t)*\ez}) -- cycle;
  \draw[blue!40] plot[domain=0:360,samples=90,variable=\t]
      ({2.3*(cos(\t)*\nx+sin(\t)*\ex)},{2.3*(cos(\t)*\ny+sin(\t)*\ey)},{2.3*sin(\t)*\ez}) -- cycle;

  % ---- inertial axes
  \draw[->,gray!70!black] (0,0,0) -- (\L,0,0) node[anchor=north east] {$X$};
  \draw[->,gray!70!black] (0,0,0) -- (0,\L,0) node[anchor=west] {$Y$};
  \draw[->,gray!70!black] (0,0,0) -- (0,0,\L) node[anchor=south] {$Z$};

  % ---- line of nodes (both directions, arrow on the positive one)
  \draw[orange!85!black,thick] (0,0,0) -- ({-2.2*\nx},{-2.2*\ny},0);
  \draw[->,orange!85!black,thick] (0,0,0) -- ({3.1*\nx},{3.1*\ny},0) node[anchor=north] {节线 $ON$};

  % ---- body axes
  \draw[->,blue!75!black,thick] (0,0,0) -- ({\L*\ax},{\L*\ay},{\L*\az}) node[anchor=west] {$x_1$};
  \draw[->,blue!75!black,thick] (0,0,0) -- ({\L*\bx},{\L*\by},{\L*\bz}) node[anchor=south west] {$x_2$};
  \draw[->,blue!75!black,thick] (0,0,0) -- ({\L*\cx},{\L*\cy},{\L*\cz}) node[anchor=south] {$x_3$};

  % ---- angle arcs
  % phi: in the XY plane, from X to ON
  \draw[->,red,thick] plot[domain=0:\ph,samples=30,variable=\t] ({1.25*cos(\t)},{1.25*sin(\t)},0);
  \node[red] at ({1.55*cos(\ph/2)},{1.55*sin(\ph/2)},0) {$\varphi$};
  % theta: in the plane containing Z and x3 (perpendicular to ON), from Z to x3
  \draw[->,red,thick] plot[domain=0:\th,samples=30,variable=\t]
      ({1.3*sin(\t)*sin(\ph)},{-1.3*sin(\t)*cos(\ph)},{1.3*cos(\t)});
  \node[red] at ({1.65*sin(\th/2)*sin(\ph)},{-1.65*sin(\th/2)*cos(\ph)},{1.65*cos(\th/2)}) {$\theta$};
  % psi: in the body x1x2 plane, from ON to x1
  \draw[->,red,thick] plot[domain=0:\ps,samples=30,variable=\t]
      ({1.6*(cos(\t)*\nx+sin(\t)*\ex)},{1.6*(cos(\t)*\ny+sin(\t)*\ey)},{1.6*sin(\t)*\ez});
  \node[red] at ({1.95*(cos(\ps/2)*\nx+sin(\ps/2)*\ex)},{1.95*(cos(\ps/2)*\ny+sin(\ps/2)*\ey)},{1.95*sin(\ps/2)*\ez}) {$\psi$};

  \fill (0,0,0) circle (1.6pt);
  \node[anchor=south east] at (0,0,0) {$O$};
  % ---- legend (screen coordinates)
  \begin{scope}[tdplot_screen_coords]
    \node[anchor=west,align=left,font=\small] at (3.3,1.2)
      {$\varphi$：$X$ 轴 $\to$ 节线（绕 $Z$ 轴）\\[2pt]
       $\theta$：$Z$ 轴 $\to$ $x_3$ 轴（绕节线）\\[2pt]
       $\psi$：节线 $\to$ $x_1$ 轴（绕 $x_3$ 轴）\\[5pt]
       $R=R_z(\psi)R_x(\theta)R_z(\varphi)$};
  \end{scope}
\end{tikzpicture}
\end{document}
